\documentclass[11pt]{article}
\usepackage[letterpaper,margin=1in]{geometry}
\usepackage{amsmath,amssymb,amsthm,mathtools,mathrsfs}
\usepackage{booktabs,tabularx,array,microtype,xurl}
\usepackage[colorlinks=true,linkcolor=blue,citecolor=blue,urlcolor=blue]{hyperref}
\newtheorem{theorem}{Theorem}[section]
\newtheorem{proposition}[theorem]{Proposition}
\newtheorem{lemma}[theorem]{Lemma}
\newtheorem{corollary}[theorem]{Corollary}
\theoremstyle{remark}
\newtheorem{remark}[theorem]{Remark}
\newcommand{\R}{\mathbb R}
\newcommand{\C}{\mathbb C}
\newcommand{\HH}{\mathcal H}
\newcommand{\KK}{\mathcal K}
\newcommand{\Aprep}{\mathcal A}
\newcommand{\Mix}{\mathfrak M_{S;p,q}}
\newcommand{\Gform}{\mathcal G}
\newcommand{\Rmet}{\mathcal R_{\mathrm{met}}}
\newcommand{\norm}[1]{\left\lVert#1\right\rVert}
\newcommand{\ip}[2]{\left\langle#1,#2\right\rangle}
\newcommand{\dd}{\,\mathrm d}
\DeclareMathOperator{\Tr}{Tr}
\DeclareMathOperator{\Ran}{Ran}
\DeclareMathOperator{\supp}{supp}
\DeclareMathOperator{\diag}{diag}
\numberwithin{equation}{section}
\title{Mixed prime response of finite Sonin projections\\
\large Exact identities and signed metric corrections}
\author{Drafted for Edward Baker}
\date{17 October 2026\\\small Unreviewed working draft}
\begin{document}
\maketitle
\begin{abstract}
We study the mixed finite difference of the positive source trace associated
with an actual finite-place Sonin projection. For a fixed compact smooth
prepared source, we retain the complete compressed inverse metric and derive
an exact mixed-response identity. The explicit prime-power bulk cancels,
so the complete response is a mixed boundary correction. It extends to a
compact selfadjoint operator on each fixed source interval, although its
four individual traces need not extend to every square-integrable source.
We isolate a finite signed Weil--Gram expression together with explicit
arithmetic, boundary and metric corrections. Whenever the source interval
is longer than \( |\log(p/q)| \), the metric correction alone is unbounded
with both signs relative to the source \(L^2\) norm. Prepared modulation
also excludes a literal identification with the full Weil form of a
nonzero fixed finite translate combination. Finally, we give an operator
approximation construction for signed bounds on the complete mixed
response, including the infinite-dimensional source complement. No
numerical sign margin or support-adapted Weil positivity estimate is obtained.
\end{abstract}

\noindent\textbf{Draft status.}
This manuscript summarizes round-one Notes 01--02 of investigation 04
\cite{note01,note02}. The round remains active. Proofs are given relative
to the explicit inherited Sonin and finite Euler inputs in
Section~\ref{sec:setup}. Internal derivations and checks used
\texttt{gpt-6.1-sol}, reasoning effort \texttt{ultra}. Cross-model
reconciliation and specialist human review remain pending.

\section{Sources, projections and inherited inputs}\label{sec:setup}

The finite-place Sonin framework transports the archimedean Sonin space
by bounded invertible Euler-factor operators. The changing metric is
essential: a Hilbertian isomorphism does not identify the ordinary
orthogonal projections without an inverse compressed metric.
Connes--Consani--Moscovici \cite[Theorem 4.6 and Section 4.8]{ccm}
provide the semilocal stability framework. We use the normalization,
smoothing and finite Euler comparison established in the repository
audits \cite{canonical,euler}; these are explicit inherited inputs,
rather than new results of the present calculation.

Work on \(\HH=L^2(\R,\dd x)\), with inner products antilinear in the first
slot, and use
\[
 \widehat F(t)=\int_\R F(x)e^{-itx}\dd x,\qquad
 U_av(x)=v(x-a),\qquad D=-i\partial_x .
\]
The inverse Fourier measure is \(\dd t/(2\pi)\). Write \(C_F\) for
ordinary convolution by \(F\), and
\begin{equation}\label{eq:kappa}
 \kappa_{F,G}(a)=\int_\R\overline{F(x)}G(x+a)\dd x,\qquad
 \kappa_F=\kappa_{F,F}.
\end{equation}
Fix \(I_L=(-L/2,L/2)\) and the prepared class
\begin{equation}\label{eq:prep}
 F=\Aprep h,\qquad \Aprep=-\partial_x^2+\tfrac14,\qquad
 h\in C_c^\infty(I_L;\C).
\end{equation}
It is exactly the compact smooth class satisfying
\(\int F(x)e^{x/2}\dd x=\int F(x)e^{-x/2}\dd x=0\), without enlarging
support \cite[Section 2]{canonical}. These conditions remove the two pole
terms. They impose neither zero mean nor additive parity.

For a finite prime set \(A\), define
\begin{equation}\label{eq:transport}
 \mathscr D_A=\prod_{j\in A}(I-j^{-1/2}U_{\log j}),\qquad
 M_A=\left.\Pi_\infty\mathscr D_A^*\mathscr D_A
                    \Pi_\infty\right|_{\KK_\infty},
 \qquad \KK_\infty=\Ran\Pi_\infty ,
\end{equation}
where \(\Pi_\infty\) is the actual cutoff-one archimedean Sonin projection
in logarithmic coordinates. The inherited orthogonal projection is
\begin{equation}\label{eq:projection}
 \Pi_A=\mathscr D_A\Pi_\infty M_A^{-1}\Pi_\infty\mathscr D_A^*,
 \qquad
 \ell_A^2I\le M_A\le u_A^2I,
\end{equation}
with \(\ell_A=\prod_{j\in A}(1-j^{-1/2})>0\) and
\(u_A=\prod_{j\in A}(1+j^{-1/2})\). The inverse acts only on
\(\KK_\infty\). All surrounding projections in \eqref{eq:projection}
implement extension to \(\HH\). The source-smoothed operators
\(C_F\Pi_A\) are Hilbert--Schmidt for compact smooth \(F\).
Put
\begin{equation}\label{eq:B}
 B_A(F,G)=\ip{C_F\Pi_A}{C_G\Pi_A}_{\mathrm{HS}},
 \qquad B_A[F]=B_A(F,F)\ge0.
\end{equation}
No trace of an unsmoothed projection is taken.

For clarity, the actual boundary package used below can be specified
without an infinite Euler product. Let \(R_+=1_{(0,\infty)}\),
\(\chi=I-R_+\), and \(\mathscr Rf(x)=f(-x)\). With
\[
 \zeta_A(s)=\prod_{j\in A}(1-j^{-s})^{-1},\qquad
 v_A(t)=\pi^{-it}
 \frac{\Gamma(1/4+it/2)}{\Gamma(1/4-it/2)}
 \frac{\zeta_A(1/2+it)}{\zeta_A(1/2-it)},
\]
the finite Euler involution is \(\mathcal F_A=\mathscr Rv_A(D)\).
Set
\begin{equation}\label{eq:boundary}
 C_A=\chi\mathcal F_A\chi,\quad T_A=\chi\mathcal F_AR_+,\quad
 Y_A=C_A(I_\chi-C_A^2)^{-1}T_A .
\end{equation}
The inherited finite-place assertions are
\[
 C_A=C_A^*\ \hbox{compact},\qquad
 T_AT_A^*=I_\chi-C_A^2\ge g_AI_\chi,\qquad g_A>0 .
\]
The inverse in \eqref{eq:boundary} acts on \(\chi\HH\).
In particular \(Y_A:R_+\HH\to\chi\HH\) is compact. Compactness of \(C_A\)
does not assert that it is Hilbert--Schmidt.

For \(a_{F,e}(t)=\tfrac12(|\widehat F(t)|^2+|\widehat F(-t)|^2)\), let
\begin{equation}\label{eq:K}
 X_{F,e}=R_+a_{F,e}(D)\chi,\qquad
 K_A[F]=2\Re\Tr_{\chi\HH}(Y_AX_{F,e}).
\end{equation}
Here \(X_{F,e}\) is trace class. The finite Euler identity
\cite[Sections 1--5]{euler}, valid for every compact smooth source,
is
\begin{equation}\label{eq:euler}
 B_A[F]=\Gamma[F]-W_A[F]+K_A[F],
\end{equation}
where
\begin{align}
 \Gamma(F,G)&=\int_\R\gamma_\infty(t)
       \overline{\widehat F(t)}\widehat G(t)\frac{\dd t}{2\pi},
 &\gamma_\infty(t)&=\Re\psi(1/4+it/2)-\log\pi, \label{eq:gamma}\\
 W_A(F,G)&=\sum_{j\in A}\sum_{m\ge1}
    (\log j)j^{-m/2}
    \{\kappa_{F,G}(m\log j)+\kappa_{F,G}(-m\log j)\}.
 \label{eq:W}
\end{align}
Polarization defines the Hermitian boundary form \(K_A(F,G)\).
The archimedean calibration underlying this package is related to
Connes--Consani \cite{cc}; the finite Euler extension and the projection
formula are imported specifically from \cite{canonical,euler}.

If \(d_F=\operatorname{diam}(\supp F)\), then \(\kappa_F(a)=0\) for
\(|a|\ge d_F\). Active powers satisfy \(m\log j<d_F\); endpoint
correlations vanish. On prepared sources the complete Weil form is
\begin{equation}\label{eq:Q}
 Q=\Gamma-W_{\mathrm{full}}.
\end{equation}
The sum is finite on each compact difference support. The equality
\(Q=B_A-K_A\) requires \(A\) to capture every active prime.
The Sonin cutoff remains one when source support changes. All constants
and limits in this manuscript concern fixed finite place sets and fixed
source intervals.

\section{The exact mixed inverse formula}\label{sec:mixed}

Fix finite \(S\) and distinct primes \(p,q\notin S\). Use one and the same
source in all four terms of
\begin{equation}\label{eq:mix}
 \Mix[F]=B_{S\cup\{p,q\}}[F]-B_{S\cup\{p\}}[F]
          -B_{S\cup\{q\}}[F]+B_S[F].
\end{equation}
Rebase on \(P=\Pi_S\), \(\KK=\Ran P\), and \(V=C_F|_\KK\).
Define
\begin{equation}\label{eq:Ej}
 r_j=j^{-1/2},\qquad \beta_j=\frac{2r_j}{1+r_j^2}<1,\qquad
 E_j=-\frac{\beta_j}{2}(U_{\log j}+U_{\log j}^*),\quad j=p,q .
\end{equation}
The ambient \(E_p,E_q\) commute. On \(\KK\), put
\[
 a=PE_pP,\quad b=PE_qP,\quad c=PE_pE_qP,\quad
 \mathsf M=I+a+b+c,
\]
\[
 N_0=V^*V,\quad N_p=V^*E_pV,\quad
 N_q=V^*E_qV,\quad N_{pq}=V^*E_pE_qV,
\]
and \(L_p=N_p-aN_0\), \(L_q=N_q-bN_0\),
\(L_{pq}=N_{pq}-cN_0\).

\begin{theorem}[Exact mixed response]\label{thm:mixed}
All \(N\)'s and \(L\)'s above are trace class, all inverses below act on
\(\KK\), and
\begin{equation}\label{eq:exactmix}
 \boxed{\Mix[F]=\Tr_\KK\mathsf M^{-1}
  \left\{L_{pq}-(b+c)(I+a)^{-1}L_p
               -(a+c)(I+b)^{-1}L_q\right\}.}
\end{equation}
Moreover,
\begin{equation}\label{eq:mixK}
 \Mix[F]=\Delta_p\Delta_qK_S[F].
\end{equation}
\end{theorem}
\begin{proof}
Hilbert--Schmidt factorization of \(V\) proves the trace-class assertions.
The metric bounds are
\[
 I+a\ge(1-\beta_p)I,\qquad I+b\ge(1-\beta_q)I,\qquad
 \mathsf M\ge(1-\beta_p)(1-\beta_q)I.
\]
The transported-space isometry, or \eqref{eq:projection} rebased on \(P\),
gives
\begin{align*}
 B_S&=\Tr N_0,\\
 B_{S\cup\{p\}}&=\Tr(I+a)^{-1}(N_0+N_p),\\
 B_{S\cup\{q\}}&=\Tr(I+b)^{-1}(N_0+N_q),\\
 B_{S\cup\{p,q\}}&=\Tr\mathsf M^{-1}(N_0+N_p+N_q+N_{pq}).
\end{align*}
For example, \(d_pd_qP(Pd_q^*d_p^*d_pd_qP)^{-1/2}\), where
\(d_j=I-r_jU_{\log j}\), is an isometry onto \(d_pd_q\KK\).
Convolution commutes with the transports. Their positive scalar
normalizations cancel between metric and numerator.

Subtract the four traces and use
\[
 \mathsf M^{-1}-(I+a)^{-1}
   =-\mathsf M^{-1}(b+c)(I+a)^{-1},
\]
and its \(p,q\) counterpart. This proves \eqref{eq:exactmix} with the
displayed factor order. Finally \(\Gamma\) is unchanged across the four
corners and \(W_A\) is additive in its prime set, so
\eqref{eq:euler} proves \eqref{eq:mixK}.
\end{proof}

If \(A_F=C_F^*C_F\), then
\begin{equation}\label{eq:returns}
 L_j=PE_j(I-P)A_FP,\quad
 L_{pq}=PE_pE_q(I-P)A_FP,\quad
 c-ab=PE_p(I-P)E_qP .
\end{equation}
Their trace-class justification is the factorization in the theorem,
not a separate assumption on \(A_FP\). Commuting ambient translations
do not make \(a,b\) commute, or remove the off-range return in
\eqref{eq:returns}. There is no mixed composite-prime label in the
arithmetic bulk, but the full response need not vanish.

\section{A compact source operator for the combined response}\label{sec:compact}

The next result extends the first-prime source regularity
\cite{closed,compact} by using compact boundary operators directly.

\begin{lemma}[Inherited crossing bound]\label{lem:cross}
For \(F\in C_c^\infty(I_L)\),
\begin{equation}\label{eq:crossbound}
 \norm{X_{F,e}}_1\le\frac L2\norm F_2^2.
\end{equation}
\end{lemma}
\begin{proof}
For the noneven multiplier insert the source interval between the
convolution factors:
\[
 R_+C_F^*C_F\chi=R_+C_F^*1_{I_L}C_F\chi.
\]
Indeed the kernel integrates
\(\overline{F(z-x)}F(z-y)\) with \(x>0>y\), which forces
\(-L/2<z<L/2\). The two factors are Hilbert--Schmidt with squared norms
\[
 \int_{I_L}(L/2-v)|F(v)|^2\dd v,\qquad
 \int_{I_L}(L/2+v)|F(v)|^2\dd v.
\]
Their sum is \(L\norm F_2^2\). The product bound and arithmetic-geometric
mean inequality give \eqref{eq:crossbound}. Conjugating the source reverses
the frequency multiplier and interchanges the two norms, proving the
same estimate for its even average. This is the argument of
\cite[Section 2]{closed}.
\end{proof}

For a compact \(Y:R_+\HH\to\chi\HH\), define the real quadratic source form
\[
 q_Y[F]=2\Re\Tr(YX_{F,e})
\]
first for arbitrary \(F\in C_c^\infty(I_L)\), without preparation.

\begin{theorem}[Compact boundary-to-source transfer]\label{thm:transfer}
The form \(q_Y\) extends to a compact selfadjoint operator \(H_{Y,L}\)
on \(L^2(I_L)\), with
\begin{equation}\label{eq:transfer}
 \norm{H_{Y,L}}\le L\norm Y,\qquad
 \norm{H_{Y,L}-H_{Y_N,L}}\le L\norm{Y-Y_N}.
\end{equation}
In particular, for
\begin{equation}\label{eq:Ymix}
 Y_{\mathrm{mix}}=Y_{S\cup\{p,q\}}-Y_{S\cup\{p\}}
                  -Y_{S\cup\{q\}}+Y_S ,
\end{equation}
the operator \(H_{\mathrm{mix},L}=H_{Y_{\mathrm{mix}},L}\) represents
\(\Mix\) on the compact smooth core.
\end{theorem}
\begin{proof}
Lemma~\ref{lem:cross} bounds \(|q_Y[F]|\) by
\(L\norm Y\norm F_2^2\). The source form satisfies the parallelogram
identity and \(q_Y[zF]=|z|^2q_Y[F]\); complex polarization and density
give a bounded selfadjoint source operator and the norm inequalities.
The map \(Y\mapsto H_{Y,L}\) is real-linear.

Approximate the full compact \(Y\) in operator norm by finite-rank \(Y_N\).
The crossing kernel is zero outside \(x,z>0,\ x+z<L\), so only the
restriction from \((0,L)\) to \((-L,0)\) contributes. Reflect the negative
output and write
\[
 \mathscr R1_{(-L,0)}Y_N1_{(0,L)}
       =\sum_j|v_j\rangle\langle u_j|,\qquad u_j,v_j\in L^2(0,L).
\]
With zero extensions define
\begin{equation}\label{eq:wkernel}
 w_N(s)=\sum_j\int_0^s\overline{u_j(x)}v_j(s-x)\dd x,\qquad
 k_N(s)=\Re w_N(|s|),\quad |s|\le L .
\end{equation}
Convolution of two \(L^2\) functions is continuous here. Thus \(k_N\)
is real, even and continuous. Direct integration gives
\begin{equation}\label{eq:sourcekernel}
 q_{Y_N}[F]=\int_{-L}^L\kappa_F(s)k_N(s)\dd s
  =\ip F{H_{Y_N,L}F},\qquad
 (H_{Y_N,L}F)(x)=\int_{I_L}k_N(y-x)F(y)\dd y .
\end{equation}
There is no additional factor two in the even kernel. This operator is
Hilbert--Schmidt, although generally not finite rank. The second bound
in \eqref{eq:transfer} makes \(H_{Y,L}\) a norm limit of compact operators.
For \eqref{eq:Ymix}, \eqref{eq:mixK} identifies the form with \(\Mix\).
\end{proof}

The compact extension concerns the combined source form. It does not
make each of the four \(B_A[F]\) traces finite on every \(L^2\) source,
or authorize subtraction of infinite traces. Nor does the proof assert
that a general multiple-prime source operator is Hilbert--Schmidt. It
uses no claim that the first-return kernel has finitely many resonances:
ratios of integers supported on two primes can be dense.

\section{A corrected signed Weil--Gram identity}\label{sec:gram}

Let
\[
 J_j=\frac{I-r_jU_{\log j}}{\sqrt{1+r_j^2}},\qquad
 J_A=\prod_{j\in A}J_j,\qquad A\subseteq\{p,q\},
\]
and, on the current space \(\KK\), set
\[
 \mathsf M_A=(PJ_A^*J_AP)|_\KK,\qquad
 \mathsf N_A=V^*J_A^*J_AV,\qquad
 Z_A[F]=\Tr\bigl((\mathsf M_A^{-1}-I)\mathsf N_A\bigr).
\]
The inverses are bounded by the fixed-place metric bounds. Put
\begin{equation}\label{eq:Rmet}
 \Rmet=Z_{\{p,q\}}-Z_{\{p\}}-Z_{\{q\}}.
\end{equation}
For \(a_+=\log(pq)\), \(a_-=|\log(p/q)|\), and any translation-invariant
Hermitian source form \(T\), define
\begin{align}
 \Gform_T[F]
 &=\frac{\beta_p\beta_q}{2}
      \Re\{T(F,U_{a_+}F)+T(F,U_{a_-}F)\} \label{eq:G}\\
 &=\frac{\beta_p\beta_q}{8}
       \sum_{a\in\{a_+,a_-\}}\{T[F+U_aF]-T[F-U_aF]\}.
 \nonumber
\end{align}
All forms used here are translation invariant. For \(B_S\), this follows
from \(C_{U_aF}=U_aC_F\) and unitary invariance of the Hilbert--Schmidt
inner product; it does not require \(\Pi_S\) to commute with translations.

\begin{theorem}[Corrected signed Gram identity]\label{thm:gram}
On the prepared smooth core,
\begin{equation}\label{eq:corrected}
 \boxed{\Mix=\Gform_Q+\Gform_{W_{\notin S}}
                         +\Gform_{K_S}+\Rmet,}
 \qquad W_{\notin S}=W_{\mathrm{full}}-W_S .
\end{equation}
Each correction has an independent arithmetic or operator definition.
Moreover,
\begin{equation}\label{eq:principal}
 \boxed{\Rmet=-\Gform_\Gamma+\Gform_{W_S}
                      +\bigl(\Delta_p\Delta_qK_S-\Gform_{K_S}\bigr).}
\end{equation}
On fixed \(L^2(I_L)\), the last parenthesis is compact and
\(\Gform_{W_S}\) is bounded.
\end{theorem}
\begin{proof}
The normalized transport formula gives
\[
 B_{S\cup A}[F]=\Tr(\mathsf M_A^{-1}\mathsf N_A)
                    =B_S[J_AF]+Z_A[F].
\]
Since \(J_j^*J_j=I+E_j\), the mixed difference of the unweighted
numerators is \(V^*E_pE_qV\). Expanding the four ambient shifts and
grouping opposite pairs gives
\(\Tr(V^*E_pE_qV)=\Gform_{B_S}[F]\), with exactly the coefficient in
\eqref{eq:G}. Hence \(\Mix=\Gform_{B_S}+\Rmet\).
Polarize \(B_S=Q+W_{\notin S}+K_S\) to obtain
\eqref{eq:corrected}. Substitute \eqref{eq:mixK} and
\(B_S=\Gamma-W_S+K_S\) to obtain \eqref{eq:principal}.

Theorem~\ref{thm:transfer} applies to every fixed \(K_A\).
The sources \(F\pm U_aF\) have a common support hull of length at most
\(L+a_+\); pullback of the corresponding compact source forms proves
compactness of \(\Gform_{K_S}\). The fixed-\(S\) prime series is a bounded
translation form, proving the remaining assertion.
\end{proof}

The complete arithmetic in \eqref{eq:corrected} is evaluated on the
enlarged translate support, so it can include primes inactive on the
original interval. On the ordered three-source Gram matrix
\(F,U_{a_+}F,U_{a_-}F\), the contraction in \eqref{eq:G} has coefficient
matrix
\[
 \begin{pmatrix}0&w&w\\w&0&0\\w&0&0\end{pmatrix},
 \qquad w=\frac{\beta_p\beta_q}{4},
\]
with eigenvalues \(0,\pm\beta_p\beta_q/(2\sqrt2)\). Positive semidefiniteness
of a Weil Gram matrix alone therefore does not sign this contraction.

\section{The metric correction has both unbounded signs}\label{sec:sign}

\begin{theorem}[Necessary signed gamma cancellation]\label{thm:sign}
Suppose \(L>a_-\). There is a prepared family \(F_\tau\) supported in
\(I_L\), with \(\norm{F_\tau}_2\to\norm h_2>0\), such that
\begin{equation}\label{eq:signasym}
 \Rmet[F_\tau]=-c_h\cos(\tau a_-)\log\tau+O(1),\qquad
 c_h=\frac{\beta_p\beta_q}{2}\kappa_h(a_-)>0 .
\end{equation}
Thus \(\Rmet[F_\tau]\to-\infty\) along \(\tau=2\pi n/a_-\), and
\(\Rmet[F_\tau]\to+\infty\) along \(\tau=(2n+1)\pi/a_-\).
The full response satisfies \(\Mix[F_\tau]\to0\).
\end{theorem}
\begin{proof}
Choose a real nonnegative smooth bump \(h\), positive on the interior of
one interval of diameter \(d\), where
\(a_-<d<\min(L,a_+)\). Then \(\kappa_h(a_-)>0\) and
\(\kappa_h(a_+)=0\). Define
\begin{equation}\label{eq:Ftau}
 F_\tau=\tau^{-2}\Aprep(e^{i\tau x}h)
       =e^{i\tau x}g_\tau,\qquad
 g_\tau=h-\frac{2i}{\tau}h'+\frac{\Aprep h}{\tau^2}.
\end{equation}
Preparation is exact, support is unchanged, and \(g_\tau\to h\) in
\(C_c^\infty\). In particular \(F_\tau\rightharpoonup0\) and its norms
tend to \(\norm h_2\).

The digamma expansion \cite[Equation 5.11.2]{dlmf} yields
\(\gamma_\infty(t)=\log|t|-\log(2\pi)+O(|t|^{-1})\).
After centering the Fourier variable at \(\tau\), uniformly Schwartz
envelopes give, for each fixed shift \(a\),
\[
 \Re\Gamma(F_\tau,U_aF_\tau)
       =(\log\tau)\kappa_h(a)\cos(\tau a)+O(1).
\]
One may split the integral into \(|t-\tau|\le\tau/2\) and its
complement. The expansion is uniform on the first region and the
Schwartz tail controls the second. Replacing \(g_\tau\) by \(h\) costs
\(O(\tau^{-1}\log\tau)\). Equation \eqref{eq:G} now gives the stated
leading term for \(\Gform_\Gamma\).

In \eqref{eq:principal}, the arithmetic contribution is bounded.
The compact terms tend to zero on this bounded weakly null family,
by Theorem~\ref{thm:transfer}. This proves \eqref{eq:signasym}.
The same theorem gives \(\Mix[F_\tau]\to0\).
\end{proof}

The theorem concerns the source quadratic form \(\Rmet\), with the
normalization \eqref{eq:Rmet}. It does not say that any compressed inverse
is an unbounded operator. It excludes a bounded \(L^2\) remainder or a
fixed sign for this metric correction, not collective arithmetic estimates
for the complete response. For \(p=2,q=3,S=\varnothing,L=6/5\), both
primes belong to the final corner, which contains every prime potentially
active on \(I_L\), and \(c_h=(\sqrt6/6)\kappa_h(\log(3/2))\). The phenomenon already occurs
for this first pair.

\section{What a finite translate identification would require}\label{sec:translate}

\begin{proposition}[Obstruction to a literal full Weil form]\label{prop:translate}
A fixed finite linear combination \(\mathcal E[F]\) of mixed responses
cannot equal
\[
 Q\left[\sum_{j=1}^n c_jU_{t_j}F\right]
\]
for every prepared \(F\in C_c^\infty(I_L)\), if the shifts are distinct
and the fixed coefficient vector \(c\) is nonzero.
\end{proposition}
\begin{proof}
Use \eqref{eq:Ftau} with any nonzero \(h\). Compactness of every fixed
mixed source operator gives \(\mathcal E[F_\tau]\to0\).
Set \(H_\tau=\sum_jc_jU_{t_j}F_\tau\). Its envelope is
\(\sum_jc_je^{-i\tau t_j}U_{t_j}g_\tau\).
Distinct translates of a nonzero compactly supported function are
linearly independent: Fourier transformation would otherwise make a
nonzero finite exponential polynomial vanish on an open set where
\(\widehat h\ne0\). Their positive definite finite Gram matrix, and
\(g_\tau\to h\), therefore give \(\norm{H_\tau}_2^2\ge\eta>0\) for
large \(\tau\), uniformly in the phases. The envelopes have uniformly
bounded Schwartz seminorms, so
\[
 \Gamma[H_\tau]=(\log\tau)\norm{H_\tau}_2^2+O(1).
\]
The complete prime sum on their fixed enlarged support stays bounded.
Preparation is preserved by translation, so \(Q[H_\tau]\to+\infty\).
\end{proof}

The proposition excludes a literal equality of the stated kind.
It does not exclude the corrected signed identity
\eqref{eq:corrected}, signed combinations of Gram entries with their
own cancellations, source-dependent coefficients, infinite
constructions, or an equivalence stated only in terms of positivity.
Likewise, any exact expression for \(\mathcal E[F]\) as a fixed linear
combination of finitely many bare autocorrelation samples must be trivial
on accessible shifts: \eqref{eq:Ftau} turns it into a finite exponential
polynomial plus \(o(1)\), while \(\mathcal E[F_\tau]\to0\).
Here a shift is accessible when \(|a|<L\); samples at \(|a|\ge L\)
vanish identically on the source class. Choose one real positive bump
whose support diameter exceeds the largest specified accessible shift
magnitude. Its autocorrelations at all those shifts are positive.
Averaging against each frequency therefore forces the grouped sample
coefficients themselves to vanish. This argument is about fixed linear
combinations, not arbitrary nonlinear functions of finitely many samples.

\section{Signed approximation bounds on the complete response}\label{sec:enclosure}

The compact source operator, rather than the unbounded metric sector
alone, is the object that admits a full-norm signed approximation.
The following construction supplies a verification contract. Its constants
have not been evaluated for an actual prime pair.

\begin{proposition}[Boundary resolvent approximation]\label{prop:resolvent}
Suppose \(C=C^*\), \(TT^*=I-C^2\ge gI\), \(\norm C\le c<1\), and
a selfadjoint finite-rank \(C_N\) obeys
\(\norm{C-C_N}\le\varepsilon\), \(\norm{C_N}\le c_N<1\).
For \(Y=C(I-C^2)^{-1}T\) and
\(Y_N=C_N(I-C_N^2)^{-1}T\),
\begin{equation}\label{eq:reserror}
 \norm{Y-Y_N}\le
 \frac{\varepsilon(1+cc_N)}{(1-c_N^2)\sqrt g}.
\end{equation}
\end{proposition}
\begin{proof}
Write \(A=I-C^2\), \(A_N=I-C_N^2\). Then
\[
 Y-Y_N=(C-C_N)A^{-1}T+
              C_NA_N^{-1}(C^2-C_N^2)A^{-1}T .
\]
Use \(\norm{A^{-1}T}\le g^{-1/2}\) and
\(\norm{C^2-C_N^2}\le(c+c_N)\varepsilon\), and simplify.
No commutation of \(C,C_N\) is assumed.
\end{proof}

Apply this proposition at the four corners, combine the finite-rank
approximants with their signs, and include any errors in the finitely
many adjoint \(T_A\) columns in a full bound
\begin{equation}\label{eq:Yerror}
 \norm{Y_{\mathrm{mix}}-Y_N}\le\varepsilon_Y .
\end{equation}
The kernel \(k_N\) in \eqref{eq:wkernel} is continuous. Divide \(I_L\)
into cells of width \(h=L/n\), with centers \(x_i\), and approximate its
source kernel on cell \(i\) times cell \(j\) by \(k_N(x_j-x_i)\).
The resulting selfadjoint \(J_{N,n}\) has rank at most \(n\). Define
\[
 \omega_N(h)=\sup\{|k_N(s)-k_N(t)|:
            s,t\in[-L,L],\ |s-t|\le h\}.
\]
The Schur bound and \eqref{eq:transfer} imply
\begin{equation}\label{eq:eta}
 \norm{H_{\mathrm{mix},L}-J_{N,n}}\le\eta,\qquad
 \eta=L\{\varepsilon_Y+\omega_N(h)\},
\end{equation}
with function evaluation and rounding errors also included.

For the prepared class compress to the closed moment kernel
\[
 \HH_L^0=\left\{F\in L^2(I_L):
       \int_{I_L}e^{x/2}F(x)\dd x
       =\int_{I_L}e^{-x/2}F(x)\dd x=0\right\}.
\]
Compression is contractive. Let \(\lambda_-,\lambda_+\) be the extremal
eigenvalues of the compressed finite-rank approximation on the span of
the projected cell indicators, computed in an orthonormal basis or with
their Gram matrix. Then
\begin{equation}\label{eq:signed}
 \boxed{
 \begin{aligned}
  \bigl(\min(0,\lambda_-)-\eta\bigr)\norm F_2^2
       &\le\Mix[F]\\
       &\le\bigl(\max(0,\lambda_+)+\eta\bigr)\norm F_2^2 .
 \end{aligned}}
\end{equation}
If the finite span is zero, take both finite spectral endpoints as zero.
The zero eigenvalue and the \(\eta I\) allowance account for the entire
infinite-dimensional source complement. A favorable finite matrix
alone cannot remove them. This is an approximation theorem, not an
already evaluated numerical certificate.

\section{Scope of the results and the remaining estimate}\label{sec:scope}

The results separate three objects: the complete mixed response, its
finite signed Weil--Gram component, and its inverse-metric correction.
Only the first is asserted to be a compact source form on every fixed
interval. The metric correction cancels an unbounded signed gamma
component whenever \(L>|\log(p/q)|\). Taking absolute values or assigning
signs before preserving that cancellation loses the relevant comparison.

Generic projection positivity cannot sign the mixed response. For example,
on \(\C^2\), let \(P_0\) project onto \((1,1)/\sqrt2\), use
\(U_a^0=\diag(1,e^{-ita})\), \(t=\pi/\log2\), and source energies
\(\diag(k_1,k_2)\), \(k_i\ge0\), obtained from
\(C_F^0=\int F(a)U_a^0\dd a\). With the actual weights at \(2,3\), put
\[
 h_{j1}=(1-r_j)^2,\quad
 h_{j2}=1+r_j^2-2r_j\cos(t\log j),\quad
 z_j=\frac{h_{j2}-h_{j1}}{h_{j2}+h_{j1}}>0 .
\]
Direct rank-one compression gives
\[
 \mathfrak M_{\varnothing;2,3}
 =-\frac{k_2-k_1}{2}
        \frac{z_2z_3(z_2+z_3)}{1+z_2z_3},
\]
which has either sign. Compact prepared sources realize arbitrary two
Fourier values at \(0,t\): the corresponding two functionals on
\(C_c^\infty(I_L)\) are independent, and preparation multiplies their
values by \(1/4,t^2+1/4\). The missing premise in this control is the
actual Sonin projection, not source neutrality or prime weights.
This is not an actual-Sonin counterexample.

The repository control screen \cite{controls} keeps further exclusions
scoped. Smooth-density controls change the discrete prime measure;
two-channel exclusions concern a specified separated architecture.
Euler-factor collision families change the arithmetic family,
weighted-square controls lack full reflection, reflected controls lack
the integer-square source, and smooth or positive-primitive controls
lack the full theta arithmetic. No noncollision theorem based on their
shared properties is claimed here.

The main outputs and their status can be stated compactly:
\begin{center}
\small
\begin{tabularx}{\textwidth}{@{}p{0.31\textwidth}X@{}}
\toprule
Output & Scope and status\\
\midrule
Mixed inverse identity & Exact on the compact smooth core, relative to
the inherited finite Sonin package.\\
Signed Weil--Gram identity & Exact with every displayed arithmetic,
boundary and metric correction retained.\\
Metric sign asymptotic & Both unbounded signs when \(L>|\log(p/q)|\);
not a sign theorem for the complete response.\\
Compact source transfer & Every fixed finite quartet and support;
no growing-place uniformity or general Hilbert--Schmidt kernel claim.\\
Signed enclosure construction & Full complement included; actual
approximation constants and numerical endpoints remain unevaluated.\\
\bottomrule
\end{tabularx}
\end{center}

The support-adapted inequality \(K_A\le B_A\), after all active primes
are captured, remains the original Weil positivity target. The present
identities do not improve its margin. The next quantitative task would
be a declared quartet and support, a useful residual budget, and certified
bounds for \eqref{eq:Yerror}--\eqref{eq:eta}. The first missing numerical
input is a full operator approximation, not a prime-pair sweep.

\section*{Acknowledgement and research record}

This working draft was prepared for Edward Baker with substantial
assistance from OpenAI Codex, recorded model \texttt{gpt-6.1-sol},
reasoning effort \texttt{ultra}. Parallel same-model agents checked the
operator algebra, source domains, modulation argument, approximation
constants and inherited source claims. These checks are internal LLM
assistance and do not constitute independent cross-model or human
refereeing. The detailed derivations are preserved in round-one Notes
01--02 \cite{note01,note02}; this manuscript is their synthesis, with no
additional numerical run. Draft history is maintained in the investigation's
\texttt{DRAFT\_HISTORY.md}, without manuscript snapshot folders.

\begin{thebibliography}{99}

\bibitem{cc}
A. Connes and C. Consani,
\emph{Weil positivity and trace formula, the archimedean place},
Selecta Math. (N.S.) \textbf{27} (2021), article 77.
\href{https://doi.org/10.1007/s00029-021-00689-4}{doi:10.1007/s00029-021-00689-4}.

\bibitem{ccm}
A. Connes, C. Consani and H. Moscovici,
\emph{Zeta zeros and prolate wave operators},
Ann. Funct. Anal. \textbf{15} (2024), article 87.
\href{https://arxiv.org/html/2310.18423v2}{arXiv:2310.18423v2};
Theorem 4.6 and Sections 4.6--4.8.

\bibitem{dlmf}
NIST Digital Library of Mathematical Functions,
Section 5.11, Equation 5.11.2.
\url{https://dlmf.nist.gov/5.11.E2}.

\bibitem{canonical}
Research note drafted for E. Baker with LLM assistance,
\emph{Sonin canonical comparison: normalization, source class, and trace audit},
29 September 2026.
Repository source:
\path{papers/weil-positivity/susy-positivity/investigations/wilson-loewner/arithmetic-storage/notes/SONIN_CANONICAL_COMPARISON_AUDIT_20260929.md}.

\bibitem{euler}
Research note drafted for E. Baker with LLM assistance,
\emph{Finite Euler boundary identity and a fixed-place sign obstruction},
3 October 2026.
Repository source:
\path{papers/investigations/sonin-critical-boundary/notes/FINITE_EULER_BOUNDARY_IDENTITY_20261003.md}.

\bibitem{closed}
Research note drafted for E. Baker with LLM assistance,
\emph{Closed source form, compact relative correction, and a finite comparison certificate},
3 October 2026, especially Section 2.
Repository source:
\path{papers/investigations/sonin-critical-boundary/notes/CLOSED_SOURCE_RELATIVE_COMPARISON_20261003.md}.

\bibitem{compact}
Research note drafted for E. Baker with LLM assistance,
\emph{Compact correction and a revised positive main term},
3 October 2026.
Repository source:
\path{papers/investigations/sonin-critical-boundary/notes/COMPACT_CORRECTION_AND_REVISED_B_20261003.md}.

\bibitem{note01}
Research note drafted for E. Baker with LLM assistance,
\emph{Exact mixed prime response and the finite translate obstruction},
17 October 2026, investigation 04, round-one Note 01.
Repository source:
\path{papers/investigations-3/04_sonin_mixed_prime_response/notes/round-1/01_EXACT_MIXED_RESPONSE_AND_FINITE_TRANSLATE_OBSTRUCTION_CODEX_20261017.md}.

\bibitem{note02}
Research note drafted for E. Baker with LLM assistance,
\emph{Corrected signed Weil Gram identity and mixed source estimates},
17 October 2026, investigation 04, round-one Note 02.
Repository source:
\path{papers/investigations-3/04_sonin_mixed_prime_response/notes/round-1/02_CORRECTED_SIGNED_WEIL_GRAM_AND_MIXED_SOURCE_ESTIMATES_CODEX_20261017.md}.

\bibitem{controls}
Repository research guide,
\emph{Collision controls and the screening they impose},
\path{result-summaries/12_collision_controls_and_screening.md}.
For inherited Sonin results and claim discipline see also
\path{result-summaries/07_sonin_critical_boundary.md},
\path{avoid/01_restating_the_target.md},
\path{avoid/07_scoped_exclusions_read_as_universal.md},
and \path{avoid/12_term_by_term_matching_and_separate_bounds.md}.

\end{thebibliography}
\end{document}
